\section*{Hoofdstuk \ref{ch:g12:stereochemistry} --- Stereochemie: chiraliteit en isomeren}

\begin{solution}{exo:g12:stereochemistry:1}
Butaan-2-ol: koolstofatoom 2 (\ce{H}, \ce{OH}, \ce{CH3}, \ce{CH2CH3}).
Propaan-2-ol: geen (koolstofatoom 2 draagt twee gelijke \ce{CH3}).
2-Chloorbutaan: koolstofatoom 2. Melkzuur: koolstofatoom 2 (\ce{H},
\ce{OH}, \ce{CH3}, \ce{COOH}).
\end{solution}

\begin{solution}{exo:g12:stereochemistry:2}
Eerste: de twee \ce{Cl} (voorrang op \ce{H}) aan dezelfde kant: \omterm{def:g12:stereochemistry:z-e}{Z}.
Tweede: links heeft \ce{CH3} voorrang, aan de bovenkant; rechts heeft
\ce{CH2CH3} voorrang, aan de onderkant: \omterm{def:g12:stereochemistry:z-e}{E}.
\end{solution}

\begin{solution}{exo:g12:stereochemistry:3}
\omterm{def:g12:stereochemistry:chiral}{Chiraal}: de handschoen, de schroef. Niet \omterm{def:g12:stereochemistry:chiral}{chiraal}: de lepel, de kubus
(elk heeft een spiegelvlak). \ce{CHFClBr} is \omterm{def:g12:stereochemistry:chiral}{chiraal} (vier verschillende
\omterm{def:g7:atoms-and-molecules:atom}{atomen} aan het koolstofatoom); \ce{CH2Cl2} niet.
\end{solution}

\begin{solution}{exo:g12:stereochemistry:4}
Een \omterm{def:g6:pure-substances-and-mixtures:mixture}{mengsel} van gelijke hoeveelheden van twee \omterm{def:g12:stereochemistry:enantiomer}{enantiomeren}. Racemisch
carvon bevat beide \omterm{def:g12:stereochemistry:enantiomer}{enantiomeren}, dus de neus krijgt beide geuren
tegelijk binnen.
\end{solution}

\begin{solution}{exo:g12:stereochemistry:5}
Nee: koolstofatoom 1 draagt twee gelijke \omterm{def:g7:atoms-and-molecules:atom}{atomen}, \ce{H} en \ce{H}; als
je die verwisselt, krijg je hetzelfde \omterm{def:g7:atoms-and-molecules:molecule}{molecuul}.
\end{solution}

\begin{solution}{exo:g12:stereochemistry:6}
\setchemfig{atom sep=2em}
\chemfig{C(-[2]CH_2CH_3)(-[4]H_3C)(<[:-30]OH)(<:[:-150]H)} \quad spiegel
\quad \chemfig{C(-[2]CH_2CH_3)(-[0]CH_3)(<[:-150]HO)(<:[:-30]H)}
\end{solution}

\begin{solution}{exo:g12:stereochemistry:7}
2-Chloorbutaan: 2 (één \omterm{def:g12:stereochemistry:chiral}{asymmetrisch koolstofatoom}). 2,3-Dichloorbutaan: 3
(twee \omterm{def:g12:stereochemistry:chiral}{asymmetrische koolstofatomen} met gelijke helften: een paar
\omterm{def:g12:stereochemistry:enantiomer}{enantiomeren} en een meso-vorm). Pent-2-een: 2 (\omterm{def:g12:stereochemistry:z-e}{Z} en \omterm{def:g12:stereochemistry:z-e}{E}).
\end{solution}

\begin{solution}{exo:g12:stereochemistry:8}
\omterm{def:g12:stereochemistry:enantiomer}{Enantiomeren}; \omterm{def:g12:stereochemistry:diastereomer}{diastereomeren}; \omterm{def:g12:stereochemistry:diastereomer}{diastereomeren}; hetzelfde \omterm{def:g7:atoms-and-molecules:molecule}{molecuul}
(meso-wijnsteenzuur valt samen met zijn spiegelbeeld).
\end{solution}

\begin{solution}{exo:g12:stereochemistry:9}
In de gestaffelde \omterm{def:g12:stereochemistry:conformation}{conformatie} met tegenover elkaar staande
\ce{CH3}-groepen zitten de grote groepen zo ver mogelijk uit elkaar: die
is het stabielst. In de eclipsconformatie staan de twee \ce{CH3} naar
elkaar toe en zitten ze elkaar in de weg.
\end{solution}

\begin{solution}{exo:g12:stereochemistry:10}
Zijn twee helften, elk \ce{CH(OH)COOH}, zijn elkaars spiegelbeeld: een
vlak tussen de twee middelste koolstofatomen snijdt het \omterm{def:g7:atoms-and-molecules:molecule}{molecuul} in twee
spiegelbeeldige helften. Een \omterm{def:g7:atoms-and-molecules:molecule}{molecuul} met een spiegelvlak is zijn eigen
spiegelbeeld: niet \omterm{def:g12:stereochemistry:chiral}{chiraal}.
\end{solution}

\begin{solution}{exo:g12:stereochemistry:11}
\setchemfig{atom sep=1.7em}
(\omterm{def:g12:stereochemistry:z-e}{Z})-pent-2-een \chemfig{C(-[:150]H_3C)(-[:210]H)=C(-[:30]CH_2-CH_3)(-[:-30]H)}
en (\omterm{def:g12:stereochemistry:z-e}{E})-pent-2-een \chemfig{C(-[:150]H_3C)(-[:210]H)=C(-[:30]H)(-[:-30]CH_2-CH_3)}.
\end{solution}

\begin{solution}{exo:g12:stereochemistry:12}
Zoals bij wijnsteenzuur, met \ce{Br} in plaats van \ce{OH} en \ce{CH3}
in plaats van \ce{COOH}: beide \ce{Br} vol, beide gestreept (een paar
\omterm{def:g12:stereochemistry:enantiomer}{enantiomeren}), en één vol en één gestreept, die een spiegelvlak heeft:
de meso-vorm. Drie \omterm{def:g12:stereochemistry:stereoisomer}{stereo-isomeren}.
\end{solution}

\begin{solution}{exo:g12:stereochemistry:13}
De helft. De scheikundige kan de twee \omterm{def:g12:stereochemistry:enantiomer}{enantiomeren} scheiden (zoals
Pasteur deed, met andere middelen), of alleen het nuttige \omterm{def:g12:stereochemistry:enantiomer}{enantiomeer}
maken, met een reactie die zelf \omterm{def:g12:stereochemistry:chiral}{chiraal} is (een \omterm{def:g12:stereochemistry:chiral}{chirale} katalysator, een
enzym, een \omterm{def:g12:stereochemistry:chiral}{chirale} uitgangsstof).
\end{solution}

\begin{solution}{exo:g12:stereochemistry:14}
Drie: (2E,4E), (2Z,4Z) en (2E,4Z), waarbij de laatste hetzelfde \omterm{def:g7:atoms-and-molecules:molecule}{molecuul}
is als (2Z,4E), van het andere eind gelezen.
\end{solution}

\begin{solution}{exo:g12:stereochemistry:15}
$2^3 = 8$ \omterm{def:g12:stereochemistry:stereoisomer}{stereo-isomeren}, dat zijn 4 paren \omterm{def:g12:stereochemistry:enantiomer}{enantiomeren}. Elk
\omterm{def:g12:stereochemistry:stereoisomer}{stereo-isomeer} heeft één \omterm{def:g12:stereochemistry:enantiomer}{enantiomeer} en $8 - 2 = 6$ \omterm{def:g12:stereochemistry:diastereomer}{diastereomeren}.
\end{solution}

\begin{solution}{pb:g12:stereochemistry:1}
\textbf{1.} \ce{C4H6O6}; $M = 4 \times 12.0 + 6 \times 1.0 + 6 \times 16.0
= \qty{150.0}{g/mol}$.

\textbf{2.} Twee carboxylgroepen (zuur) en twee hydroxylgroepen
(\omterm{def:g11:functional-groups:alcohol}{alcohol}).

\textbf{3.} De twee middelste koolstofatomen, elk gebonden aan \ce{H},
\ce{OH}, \ce{COOH} en \ce{CH(OH)COOH}.

\textbf{4.} $2^2 = 4$.

\textbf{5.} Allebei met volle wiggen.

\textbf{6.} Beide \ce{OH} met gestreepte wiggen; het valt niet samen met
L: het is D-wijnsteenzuur, zijn \omterm{def:g12:stereochemistry:enantiomer}{enantiomeer}.

\textbf{7.} Als je het \omterm{def:g7:atoms-and-molecules:molecule}{molecuul} om de middelste binding draait zodat de
twee \ce{OH} dezelfde kant op wijzen, heeft het een vlak tussen zijn
twee middelste koolstofatomen dat de ene helft op de andere spiegelt.

\textbf{8.} Nee: het is meso-wijnsteenzuur.

\textbf{9.} Van de vier combinaties (vol, vol), (gestreept, gestreept),
(vol, gestreept), (gestreept, vol) zijn de laatste twee hetzelfde
\omterm{def:g7:atoms-and-molecules:molecule}{molecuul}, de meso-vorm, omdat de twee helften van het \omterm{def:g7:atoms-and-molecules:molecule}{molecuul} gelijk
zijn.

\textbf{10.} L en D: \omterm{def:g12:stereochemistry:enantiomer}{enantiomeren}. L en meso: \omterm{def:g12:stereochemistry:diastereomer}{diastereomeren}.

\textbf{11.} Ook bij \qty{169}{\celsius}: \omterm{def:g12:stereochemistry:enantiomer}{enantiomeren} hebben dezelfde
fysische eigenschappen.

\textbf{12.} Nee: het is een \omterm{def:g12:stereochemistry:diastereomer}{diastereomeer} van L, een andere verbinding
met eigen eigenschappen.

\textbf{13.} De receptoren van de neus zijn \omterm{def:g12:stereochemistry:chiral}{chiraal} en onderscheiden
spiegelbeelden; een thermometer, die niet \omterm{def:g12:stereochemistry:chiral}{chiraal} is, kan dat niet.

\textbf{14.} Nee, het is een \omterm{def:g6:pure-substances-and-mixtures:mixture}{mengsel} van twee verbindingen; als geheel is
het niet optisch actief, omdat de twee \omterm{def:g12:stereochemistry:enantiomer}{enantiomeren} elkaars effect
opheffen.

\textbf{15.} Maar één \omterm{def:g12:stereochemistry:enantiomer}{enantiomeer}: elk kristal was opgebouwd uit alleen L
of alleen D.

\textbf{16.} \qty{5.0}{g} van elk.

\textbf{17.} De meso-vorm is een andere verbinding, geen deel van het
\omterm{def:g12:stereochemistry:enantiomer}{racemische mengsel} van L en D.

\textbf{18.} Door elk hoopje op te lossen: de oplossingen draaien
gepolariseerd licht over dezelfde hoek in tegengestelde richting (en de
kristallen zijn elkaars spiegelbeeld).

\textbf{19.} Vier: zijn twee \omterm{def:g12:stereochemistry:chiral}{asymmetrische koolstofatomen} dragen
verschillende groepen (\ce{Cl} aan het ene, \ce{OH} aan het andere), dus
geen enkele combinatie heeft een spiegelvlak, en geen twee combinaties
zijn hetzelfde \omterm{def:g7:atoms-and-molecules:molecule}{molecuul}.

\textbf{20.} Drie: L, D en meso.
\end{solution}
